\documentclass[spanish,letterpaper,oneside]{article}

\def\documenttitle {Trabajo Final: Una habilidad, un PDA, cinco minutos · Ficha de trabajo}
\def\documentsubtitle {Temas selectos de matemáticas I · Cohorte SD26}
\def\documentsubject {Actividad 6 · Material para estudiantes}
\def\documentauthor {Mtro. Martín Jonathan de la Cruz Muñoz}
\def\coursename {Temas selectos de matemáticas I}
\def\coursecode {PFM02}
\def\universityname {Instituto de Investigación, Innovación y Estudios de Posgrado para la Educación del Estado de Nuevo León}
\def\universityfaculty {Maestría en Aprendizaje y Enseñanza de las Matemáticas}
\def\universitydepartment {Temas selectos de matemáticas I}
\def\universitydepartmentimage {img/iiiepe/marca-iiiepe.png}
\def\universitydepartmentimagecfg {height=1.75cm}
\def\universitylocation {Monterrey, Nuevo León}
\def\studentfullname {Mtro. Martín Jonathan de la Cruz Muñoz}
\def\studentid {00874}
\def\studentemail {martin.delacruz@campus.iiiepe.edu.mx}
\def\programperiod {2do periodo, agosto--noviembre de 2026 (SD26)}
\def\programmodality {Modalidad presencial y en línea}
\def\programshort {MAEM}

\def\authortable {
  \begin{tabular}{ll}
    \textbf{Docente:} & \studentfullname \\
    Matrícula: & \studentid \\
    Unidad: & \coursename \\
    Periodo: & \programperiod \\
    \multicolumn{2}{l}{Fecha: 14 de noviembre de 2026}
  \end{tabular}
}

\input{template}

\usepackage{helvet}
\renewcommand{\familydefault}{\sfdefault}
\usepackage{amsmath,amssymb,array,longtable,tikz,tabularx,colortbl,enumitem}
\definecolor{iiiepeNavy}{HTML}{163B5C}
\definecolor{iiiepeBlue}{HTML}{1F6F9F}
\definecolor{iiiepePale}{HTML}{EEF5F8}
\def\portraittitlecolor {iiiepeNavy}

\setlength{\parskip}{5pt}\setlength{\parindent}{0pt}
\begin{document}\templatePortrait\templatePagecfg\templateFinalcfg\fancyfoot{}\fancyfoot[L]{\scriptsize IIIEPE · MAEM · SD26}\fancyfoot[C]{\scriptsize\thepage}\fancyfoot[R]{\scriptsize Ficha Act. 6}\setstretch{1.2}\raggedbottom

\section*{Reto del Trabajo Final: Sistema hidráulico proporcional}
Nombre: \rule{8cm}{.2pt}\hfill Fecha: \rule{2.5cm}{.2pt}

Un brazo hidráulico cuenta con un émbolo de entrada de diámetro $d_1 = 5\text{ cm}$ y un émbolo de salida de diámetro $d_2 = 10\text{ cm}$. Al aplicar una fuerza $F_1 = 250\text{ N}$ desplazando el émbolo menor $\Delta x_1 = 20\text{ cm}$:
\begin{enumerate}
  \item Calcula el factor cuadrático de amplificación $k = (d_2/d_1)^2$ y deduce $F_2 = k\cdot F_1$.
  \item Determina el desplazamiento resultante $\Delta x_2$ por conservación de trabajo $F_1\Delta x_1 = F_2\Delta x_2$.
  \item Completa una tabla para $F_1 \in \{0, 100, 200, 250\}\text{ N}$ y explica por qué suponer $F_2 = 2 F_1$ es erróneo.
\end{enumerate}
\vspace{3.5cm}

\end{document}
